% from ega4-9-fr.tex
\oldpage[IV-3]{94}

vérifie la propriété géométrique $(R_k)$ au point $y\in f^{-1}(s)$ signifie, ou
bien que l'on a $y\in L$, ou bien que les points génériques $z_i$ des
composantes irréductibles de l'ensemble fermé $F_s$ qui contiennent $y$
vérifient la relation $\dim(\mathcal O_{X_s,z_i})\geq k+1$ (compte tenu de
(4.2.7) et (5.2.3))~; autrement dit, les points $y\in F_s$ où $X_s$ vérifie la
propriété $(R_k)$ géométrique sont ceux tels que
$\operatorname{codim}_y(F_s,X_s)\geq k+1$ (5.1.2). La conclusion résulte donc
de (i) et de (9.9.1, (ii)).

Enfin, dans le cas (iii), la conclusion résulte de (ii), de (9.9.2, (viii)), du
fait que la propriété $(S_k)$ est stable par extension du corps de base (6.7.1)
et enfin du critère de Serre (5.8.6).

\begin{corollary}[9.9.5]
\label{IV.9.9.5-fr}
Soit $f:X\to S$ un morphisme de présentation finie, et notons $\mathbf P$
l'une des propriétés (i) à (v) de (9.9.4). Alors l'ensemble des $s\in S$ tels
que la propriété $\mathbf P$ soit vraie en tous les points $x\in X_s$, est
localement constructible dans $S$.
\end{corollary}

La démonstration à partir de (9.9.4) est la même que celle de (9.9.3) à partir
de (9.9.2).

\begin{proposition}[9.9.6]
\label{IV.9.9.6-fr}
Soient $f:X\to S$ un morphisme localement de présentation finie,
$\mathcal L_\bullet$ un complexe formé de $\mathcal O_X$-Modules
quasi-cohérents de présentation finie~; pour tout $s\in S$, soit
$(\mathcal L_\bullet)_s$ le complexe
$\mathcal L_\bullet\otimes_{\mathcal O_S}k(s)$ de
$\mathcal O_{X_s}$-Modules cohérents de type fini. Alors, pour un entier $n$
donné, l'ensemble des $x\in X$ tels que
$(\mathcal H_n((\mathcal L_\bullet)_{f(x)}))_x=0$ est localement constructible
dans $X$.
\end{proposition}

On peut se borner au cas où $\mathcal L_i=0$ sauf pour $i=0$, 1 ou 2, et où
$n=1$. En outre, la question étant locale sur $X$, on peut se borner au cas où
$S=\operatorname{Spec}(A)$ et $X=\operatorname{Spec}(B)$ sont affines, $B$
étant une $A$-algèbre de présentation finie. Il existe alors un sous-anneau
noethérien $A_0$ de $A$, un $A_0$-préschéma de type fini $X_0$ et un complexe
$\mathcal L_\bullet^{(0)}$ de $\mathcal O_{X_0}$-Modules cohérents, nuls sauf en
dimensions 0, 1 et 2 et tels que
$X=X_0\otimes_{A_0}A$ et
$\mathcal L_\bullet=\mathcal L_\bullet^{(0)}\otimes_{A_0}A$. Pour tout
$s\in S$, si $s_0$ est la projection de $s$ dans
$S_0=\operatorname{Spec}(A_0)$, on a
$X_s=(X_0)_{s_0}\otimes_{k(s_0)}k(s)$, et le morphisme projection
$X_s\to(X_0)_{s_0}$ est fidèlement plat~; on en conclut que l'on a
\[
\mathcal H_n((\mathcal L_\bullet)_s)=
\mathcal H_n((\mathcal L_\bullet^{(0)})_{s_0})
\otimes_{k(s_0)}k(s),
\]
et par suite, si $E$ (resp. $E_0$) est l'ensemble des $x\in X$ (resp.
$x_0\in X_0$) tels que
\[
(\mathcal H_n((\mathcal L_\bullet)_{f(x)}))_x=0
\quad\text{(resp.}\quad
(\mathcal H_n((\mathcal L_\bullet^{(0)})_{f_0(x_0)}))_{x_0}=0\text{)},
\]
on a $E=h^{-1}(E_0)$, où $h:X\to X_0$ est la projection canonique. En vertu
de (1.8.2), on peut donc se borner au cas où $A$ est \emph{noethérien}~; il
s'agit de voir ($\mathbf 0_{\mathrm{III}},9.2.3$) que si $x\in X$ est tel que
$(\mathcal H_n((\mathcal L_\bullet)_{f(x)}))_x=0$ (resp. $\neq0$), il existe un
voisinage ouvert $V$ de $x$ dans $\overline{\{x\}}$ tel que, pour tout
$x'\in V$, on ait
$(\mathcal H_n((\mathcal L_\bullet)_{f(x')}))_{x'}=0$ (resp. $\neq0$).
Remplaçant $S$ par le sous-préschéma réduit de $S$ ayant
$\overline{\{f(x)\}}$ pour espace sous-jacent, on peut supposer que $S$ est
intègre et que $f(x)=\eta$ est son point générique. Alors, en restreignant $S$
à un voisinage ouvert de $\eta$, on peut supposer que pour tout $s\in S$, on a
$(\mathcal H_n(\mathcal L_\bullet))_s=
\mathcal H_n((\mathcal L_\bullet)_s)$ (9.4.3), et par suite, si $Z$ est le
support de $\mathcal H_n(\mathcal L_\bullet)$, le support de
$\mathcal H_n((\mathcal L_\bullet)_s)$ est $Z_s=Z\cap X_s$
(\textbf{I}, 9.1.13.1). L'hypothèse est que
$Z_\eta\cap\overline{\{x\}}=\varnothing$ (resp. $\neq\varnothing$). Comme
$Z\cap\overline{\{x\}}$ est fermé dans l'espace noethérien $X$, on conclut de
(9.5.1) qu'il y a un voisinage $U$ de $\eta$ dans $S$ tel que, pour tout
$s\in U$, on ait $Z_s\cap\overline{\{x\}}=\varnothing$ (resp.
$\neq\varnothing$). L'ensemble $V=f^{-1}(U)\cap\overline{\{x\}}$ répond donc à
la question.
